% Build the layer, Chapter 2: a verifier for cards.
% Inserted in \S\ref{sec:architecture}, after Phase 3 and before the
% subsection on revoking a card.
%
% No code or output in this file is typed by hand. \pdexcerpt, \pdrun and
% \pdrunlines print files under excerpts/, which the kernel's
% tools/excerpts.mjs cuts from its source and its recorded runs (see
% figures/pd-listing.tex). The figure draws one of those recorded runs as it
% was printed. To change a listing or the figure, change the kernel and cut
% again. The prose names functions, never line numbers, so that a listing
% which moves in its file does not make a sentence wrong.
\subsection{Build the layer: a verifier for delegated cards}
\label{sec:anchor-build}

\noindent\emph{Express lane: this section builds a verifier twice. The first
reads a card's header to choose how to verify it, and it admits a forged card
and two chains of validly signed ones. The second verifies every card with
one algorithm and compares each card with its parent. Every chain of three
cards over two rights is then presented to both.}

Phases 1 to 3 said what a verifier must do, and this section builds one. The
code is \pdkernelname{}, the small kernel written for this book. It is not
the daemon whose verifier the rest of this chapter describes.

The section has the same four parts as \weld{prereq:sec:prereq-build}{Naive}:
\emph{Naive}, \emph{Break}, \emph{Fix} and \emph{Verify}. As there, a
script cuts each listing from the kernel's source, and each run is a
recording. The build stops where Phase~3 stops. Nothing below revokes a
card, and nothing ties a card to the process that presents it.
\S\ref{sec:anchor-build-taken-back} builds revocation, and
\S\ref{sec:anchor-build-holder} makes the presenter prove that it holds the
card's key.

\subsubsection{Naive: read the header to choose the algorithm}

\pdkernelname{} encodes a card as this chapter does, as a JSON Web Token.
Listing~\ref{alg:anchor-build-cards} decodes a chain of three.

\pdrun[caption={Three cards, decoded. Each names the key that the next must
be signed with.},label={alg:anchor-build-cards}]{cards}

Card 1 is the harbor's. Its \texttt{sub} is the holder's name, and its
\texttt{key} is the holder's public key. \texttt{can} lists the holder's
rights. A right is a verb, \texttt{read} or \texttt{write}, on a path, and a
path that ends in a star matches every path that starts with the text before
the star. \texttt{exp} is the card's expiry time, in milliseconds on a clock
that the runs below set explicitly, so 600000 is ten minutes. \texttt{jti} is
a serial number. The harbor gives one to each card it issues.

Alice made card 2 without the harbor. She wrote its claims and signed it
with her own private key, the one whose public half card 1 names. A card
made this way has no serial number. It carries \texttt{par}, the identifier
of its parent. For card 3 that identifier is the SHA-256 of card
2's signature, which is the rule of \S\ref{sec:algorithms} for a link that
no issuer numbered.

The verifier holds one key, the harbor's public key. Someone writing it
would reasonably start from the format. A token's header names the algorithm
that signed it, so read the header and verify the card with that algorithm.
For a chain, verify each signature with the key that the parent names, and
then make sure the last card grants nothing that the harbor's card did not.
Listing~\ref{alg:anchor-build-naive} is that verifier.

\Needspace{14\baselineskip}
\pdexcerpt[caption={The naive verifier. \texttt{signed} reads the header to
choose how to verify a card. \texttt{check} compares the last card with the
first.},label={alg:anchor-build-naive}]{ch02-naive}

This is careful code. \texttt{signed} knows two algorithms, as a
general-purpose library would, and refuses any other. \texttt{check} refuses
an expired card, and \texttt{grants} knows the rule for stars.
It gets the chapter's own example right
(Listing~\ref{alg:anchor-build-handing-on}).

\pdrun[caption={Delegation, under the naive verifier. Every request gets the
answer its cards require.},label={alg:anchor-build-handing-on}]{delegation-naive}

Alice's card had ten minutes left. She requested twenty minutes for Heron,
and the card she signed is valid for ten, because the code that delegates a
card never signs past the expiry of its parent. Heron delegated to Finch for
five minutes. Finch's card expired first. Heron's expired five minutes
later, and Alice's own card with it.

\subsubsection{Break: the attacker writes the header}

Wren has no card. She has the harbor's public key, as everyone does. In the
runs that follow, the \emph{gate} is the verifier.

\pdrun[caption={Algorithm confusion. The second of Wren's cards is
admitted.},label={alg:anchor-build-confusion-naive}]{confusion-naive}

Wren tried the same card with three headers. The verifier refused the
header that says \texttt{none}, and it refused the card she signed with her
own key. It admitted the middle one. That card's header says HS256, so
\texttt{signed} computed a message authentication code over the card with the
only key it holds, and compared the result with the card's tag. Wren had
computed the same code with the same key. The key is the harbor's public key,
which is published so that anyone can verify a signature. As the secret of a
message authentication code, a MAC for short, it is a secret that everyone
has.

Nothing in the cryptography failed. The header chose the algorithm, and
Wren wrote the header. This is the flaw recorded as
CVE-2015-9235~\pdcite{cve20159235}, found in several token libraries at
once~\pdcite{mclean2015critical}.

\subsubsection{Break again: the last card is compared with the first}

The second break needs no forgery. Every signature in
Listing~\ref{alg:anchor-build-escalation-naive} is valid.

\pdrun[caption={Escalation. Heron may not write the file, and the card he
signs lets Finch write it.},label={alg:anchor-build-escalation-naive}]{escalation-naive}

Alice may read and write under \texttt{shop/}. She delegated to Heron the
right to read one file, and the verifier correctly refused Heron's own
request to write it. Heron then signed Finch a card to write the file.
Heron's signature on it is valid, since he signed with the key that his card
names. Writing the file is within what card 1 grants. The card passes both
checks, and Finch may do what Heron, who signed Finch's card, may not.
Figure~\ref{fig:anchor-v6v7-escalation} draws this chain, and the ProVerif
model called v6 in \S\ref{sec:attn-soundness} makes the same mistake.

\subsubsection{Break a third time: the card is presented under another parent}

The third break is less obvious. The naive verifier takes a valid signature
as proof that a card was delegated from the card presented as its parent. A
signature shows which key signed. It does not show which card the signer
was delegating from.

\Needspace{14\baselineskip}
\pdrun[caption={A moved card. Alice delegated it from her narrow card, and
Heron presents it under her wide one.},label={alg:anchor-build-moved-naive}]{moved-naive}

Alice holds two cards from the harbor, and both name her key. The second
grants \texttt{read shop/checkout.ts} and nothing more. From it she
delegates what Heron requests, which is \texttt{write shop/checkout.ts}.
That is not careless of her. The code that delegates a card,
\texttt{delegate} in the kernel's source, caps the card's expiry and does not
inspect its rights, on the understanding that the verifier refuses a card
that grants more than its parent. The verifier does refuse it, under the
narrow card. Heron then presents Alice's wide card as the parent. Her
signature is as valid there, and writing is within the wide card, so he is
admitted. A holder may well have several cards: one for standing work, and
one narrowed for a single task.

\subsubsection{Fix: one algorithm, and each card compared with its parent}

The fix is in the kernel's decision core, the Rust function that decides
every proposal in this book's code
(\weld{prereq:sec:prereq-build}{decision core}).
Listing~\ref{alg:anchor-build-signed} is its check of one card.

\Needspace{14\baselineskip}
\pdexcerpt[caption={One card. The signature is verified before any of the
card's claims is read.},label={alg:anchor-build-signed}]{ch02-signed}

\texttt{signed} is given a key by its caller and verifies with one
algorithm, \texttt{ed25519}. Nothing in the card is read before that, and
nothing in the card can change it. Then the header is compared, whole, with
the one header the harbor writes. It is never parsed, so no field of it is
trusted. Last, the claims are parsed, and a card with an unknown field is
refused.

The check is the strict form of Ed25519. The scheme as first deployed lets
anyone turn a valid signature into a second valid one for the same message,
and libraries differ on whether they accept the
second~\pdcite{chalkias2020taming}. A delegated card is identified by the
hash of its signature. A verifier that accepted both forms would have two
identifiers for one card.

Listing~\ref{alg:anchor-build-admit} checks a chain.

\Needspace{14\baselineskip}
\pdexcerpt[caption={A chain. Each card is verified with the key that its
parent names.},label={alg:anchor-build-admit}]{ch02-admit}

The key for a card comes from its parent, and by then the parent's own
signature has been verified. For the first card the key is the harbor's. The
\texttt{match} separates the two kinds of card in
Listing~\ref{alg:anchor-build-cards}. A first card must carry a serial
number. A later card must name its parent, which is what
\texttt{names\_parent} checks, and \texttt{within} then compares it with
that parent (Listing~\ref{alg:anchor-build-within}). Anything else is
refused. A card's expiry is checked wherever the card is in the chain, and
at the end the last card must grant the requested right.

\Needspace{14\baselineskip}
\pdexcerpt[caption={The order on grants, for rights written as a verb on a
path.},label={alg:anchor-build-within}]{ch02-within}

\texttt{within} is the order of \S\ref{sec:architecture}. Each right on the
card must be covered by a right on its parent, and the card must expire no
later. \texttt{includes} is the whole rule for stars. The naive verifier had
the same rule and the same loop over rights. What differs is the card that
\texttt{within} is given. Here it is the parent. There it was the first
card.

% A page may break here. Without this the two last lines above are carried
% over with the key idea, which is set to stay with the text before it.
\pagebreak[0]
\keyidea[The verifier decides what a card is checked against]{A card is the
input being checked. If its header chooses the algorithm, its author has
chosen the algorithm. If it is compared with the harbor's card, or with
whichever card is presented as its parent, then what its signer actually
held is not checked. In each break a question that belongs to the verifier
was answered by the party being checked.}

\subsubsection{Verify: the same runs, then every chain}

The honest chain first. The fixed verifier gives every request of
Listing~\ref{alg:anchor-build-handing-on} the answer the naive one gave,
and each refusal now names the card that caused it
(Listing~\ref{alg:anchor-build-handing-on-core}).

\Needspace{14\baselineskip}
\pdrunlines[caption={Delegation, against the fix, after time has
passed.},label={alg:anchor-build-handing-on-core}]{delegation}{1,18-23,26-27}{1, 18\textendash 23, 26 and 27}

Then Wren (Listing~\ref{alg:anchor-build-confusion-core}).

\Needspace{14\baselineskip}
\pdrunlines[caption={Algorithm confusion, against the fix. Each card fails
its signature check.},label={alg:anchor-build-confusion-core}]{confusion-core}{1,11-28}{1 and 11\textendash 28}

All of her cards are refused, and for one reason. The card whose header
says HS256 was never verified as HS256. Its tag is a MAC, and a MAC is not an
Ed25519 signature of anything. The comparison of headers was never reached.
Design invariant~\ref{prop:alg-pinning} says why it should not have to be:
the algorithm was fixed before the card arrived.

Next Finch, whose card grants what Heron's does not
(Listing~\ref{alg:anchor-build-escalation-core}).

\Needspace{13\baselineskip}
\pdrunlines[caption={Escalation, against the fix. The refusal names both
cards.},label={alg:anchor-build-escalation-core}]{escalation-core}{1,14-21}{1 and 14\textendash 21}

Last the moved card. Under Alice's narrow card it grants too much. Under her
wide card it names a different parent
(Listing~\ref{alg:anchor-build-moved-core}).

\Needspace{16\baselineskip}
\pdrunlines[caption={The moved card, against the fix. Heron's card names
Alice's narrow card, and is refused under any other.},label={alg:anchor-build-moved-core}]{moved-core}{1,11-21}{1 and 11\textendash 21}

One chain refused is one chain. Three cards that each grant some of two
rights make 64 chains. \pdkernelname{} builds every one of them with real
signatures and presents it to both verifiers
(Listing~\ref{alg:anchor-build-chains}).

\Needspace{14\baselineskip}
\pdrun[caption={Every chain of three cards, by kind, and how many of each
kind a verifier admits.},label={alg:anchor-build-chains}]{chains}

In 16 chains no card grants more than its parent, and both verifiers
accept those. The naive verifier accepts 20 more. In 11 of them the last
card has a right that the middle card lacks, which is Finch's case. In the
other 9 the middle card has a right the first lacks, and the last card
does not use it. Those 9 give the last holder nothing the harbor did not
grant. Each still has a middle card that was never a valid card.

Figure~\ref{fig:anchor-64-chains} shows where the counts come from.
Take one right. Each of the three cards has it or does not, which makes
eight combinations. In four of them no card has the right unless its parent
does. The check at every card accepts a chain when both rights are in one
of those four combinations, and $4\times4=16$. Comparing the last card with
the first allows two more combinations, and $6\times6=36$. One of the two is
the right on the first and last cards and not on the middle one. A chain
with a right placed like that is an escalation, and there are $36-25=11$ of
them.
\emph{Now you try:} Exercise~\ref{ex:anchor-build-count} asks for the 9.

% Brief (grammar: grid of cells, one for each case; atlas
%   V/fig:anchor-64-chains).
% Reader question: of all the chains a holder could build, which does each
% verifier accept, and why are the counts 16 and 36?
% Claim: comparing each card with the card above it accepts exactly the
% chains in which no card holds more than the one above. Comparing the last
% card with the first accepts 20 more, and in 11 of those the last holder
% has a right that the signer of its card lacked.
% Evidence: every chain of three cards over two rights on one file, each
% with real signatures, shown to both verifiers [internal, harbor chains
% --each; session chains-each.txt, drawn as printed].
% Must distinguish: accepted by both; accepted by the naive verifier alone,
% harmful and harmless; refused by both.
% Grammar chosen: an 8 by 8 grid with one axis for each right. A right can
% sit on three cards in eight ways, so each verifier's condition is a
% condition on each axis separately, and what it accepts is a square.
% Grammar rejected: a table of counts, which gives 16, 36, 11 and 9 and not
% the reason for them.
% Five-second test: which cells does the dashed square add to the solid one?
\begin{figure}[H]
\centering
\pdfigurehue{pdcobalt}%
\begin{tikzpicture}[pd figure,x=.6cm,y=.6cm,
    cell/.style={minimum size=.6cm-1.8pt,inner sep=0pt,outer sep=0pt},
    % Hatching drawn as strokes. A pattern fill is a tile at the page's
    % origin, which the figure check reads as ink outside the page.
    idle/.style={draw=pdinkmuted!55,line width=.5pt,path picture={
      \foreach \shift in {-.5,-.25,0,.25,.5}{
        \draw[draw=pdinkmuted!75,line width=.3pt]
          ([xshift=\shift*.6cm]path picture bounding box.south west) --
          ([xshift=\shift*.6cm]path picture bounding box.north east);}}},
    held/.style={draw=pdink,line width=.45pt,fill=pdink},
    lacking/.style={draw=pdink,line width=.45pt,fill=pdpage}]
  % The table is the recorded run, read as it was printed: the first line
  % is the command, and every remark starts with '#'.
  \pgfplotstableread[skip first n=1]{excerpts/chains-each.txt}\pdchains
  \pgfplotstablegetrowsof{\pdchains}
  \pgfmathtruncatemacro{\pdlastchain}{\pgfplotsretval-1}
  % Four counts, taken from the table as it is drawn, so that no number in
  % the figure is typed.
  \ifdefined\pdchainsboth\else
    \newcount\pdchainsboth \newcount\pdchainsgrown
    \newcount\pdchainsidle \newcount\pdchainsneither
  \fi
  \global\pdchainsboth=0 \global\pdchainsgrown=0
  \global\pdchainsidle=0 \global\pdchainsneither=0
  \foreach \i in {0,...,\pdlastchain}{
    \pgfplotstablegetelem{\i}{r}\of\pdchains \let\r\pgfplotsretval
    \pgfplotstablegetelem{\i}{w}\of\pdchains \let\w\pgfplotsretval
    \pgfplotstablegetelem{\i}{kind}\of\pdchains \let\kind\pgfplotsretval
    \pgfplotstablegetelem{\i}{naive}\of\pdchains \let\naive\pgfplotsretval
    \pgfplotstablegetelem{\i}{core}\of\pdchains \let\core\pgfplotsretval
    \ifnum\core=1
      \global\advance\pdchainsboth 1 \node[cell,pd focus fill] at (\r,-\w) {};
    \else\ifnum\naive=0
      \global\advance\pdchainsneither 1 \node[cell,pd hairline] at (\r,-\w) {};
    \else\ifnum\kind=1
      \global\advance\pdchainsgrown 1 \node[cell,pd breach fill] at (\r,-\w) {};
    \else
      \global\advance\pdchainsidle 1 \node[cell,idle] at (\r,-\w) {};
    \fi\fi\fi
  }
  % The eight ways a right can sit on three cards, in the order the run
  % prints. A filled dot is a card that holds the right.
  \foreach \bits [count=\n from 0] in {{0,0,0},{1,0,0},{1,1,0},{1,1,1},{0,1,0},{1,0,1},{0,1,1},{0,0,1}}{
    \foreach \bit [count=\k from 0] in \bits {
      \ifnum\bit=1 \def\pdplacedot{held}\else \def\pdplacedot{lacking}\fi
      \path[\pdplacedot] (\n-.28+\k*.28,1.05) circle[radius=.1];
      \path[\pdplacedot] (-1.78+\k*.28,-\n) circle[radius=.1];
    }
  }
  \node[pd label,anchor=south west] at (-.5,1.45) {the cards that hold the right to read};
  \node[pd label,rotate=90,anchor=south east] at (-2.35,.5) {the cards that hold the right to write};
  % What each verifier accepts is a square in the corner. The dashed one
  % is drawn a little outside, so both show where they share an edge.
  \draw[pd breach rule] (-.62,.62) rectangle (5.56,-5.56);
  \draw[pd focus rule] (-.5,.5) rectangle (3.5,-3.5);
  \draw[pd focus rule] (8.9,-.4) -- ++(1.1,0);
  \node[pd label,anchor=north west,align=left] at (8.8,-.6)
    {the fix compares each\\card with the card\\above it: accepts \the\pdchainsboth};
  \draw[pd breach rule] (8.9,-3.6) -- ++(1.1,0);
  \node[pd label,anchor=north west,align=left] at (8.8,-3.8)
    {the naive verifier\\compares the last card\\with the first: accepts \the\numexpr\pdchainsboth+\pdchainsgrown+\pdchainsidle\relax};
  % The key, drawn with the cells' own styles.
  \node[cell,pd focus fill] (k0) at (-.2,-8.5) {};
  \node[pd label,anchor=west] at (k0.east) {\ \the\pdchainsboth, both accept: no card holds more than the card above it};
  \node[cell,idle] (k1) at (-.2,-9.55) {};
  \node[pd label,anchor=west] at (k1.east) {\ \the\pdchainsidle, naive alone: only the middle card holds too much};
  \node[cell,pd breach fill] (k2) at (-.2,-10.6) {};
  \node[pd label,anchor=west] at (k2.east) {\ \the\pdchainsgrown, naive alone: the last card holds a right the middle lacks};
  \node[cell,pd hairline] (k3) at (-.2,-11.65) {};
  \node[pd label,anchor=west] at (k3.east) {\ \the\pdchainsneither, both refuse: the last card holds a right the first lacks};
\end{tikzpicture}
\caption{\textbf{What each verifier accepts is a square, and the naive
verifier's square is the larger.} Every chain of three cards over two
rights. The three dots on a column say which cards hold the right to read,
first to last, and those on a row say the same for writing. [internal,
\texttt{harbor chains -{}-each}]}
\label{fig:anchor-64-chains}
\end{figure}


Last, the guard. \pdkernelname{} keeps a second build of its core in which
\texttt{limit}, the function that Listing~\ref{alg:anchor-build-admit}
calls to choose what a card is compared with, returns the first card.
Finch is admitted again (Listing~\ref{alg:anchor-build-escalation-mutant}),
and the count changes (Listing~\ref{alg:anchor-build-chains-mutant}).

\Needspace{10\baselineskip}
\pdrunlines[caption={Escalation, against the fix with its guard
changed.},label={alg:anchor-build-escalation-mutant}]{escalation-compare-root}{1,16-20}{1 and 16\textendash 20}

\Needspace{9\baselineskip}
\pdrunlines[caption={The fix with its guard changed to compare with the
first card.},label={alg:anchor-build-chains-mutant}]{chains-compare-root}{1,5-8}{1 and 5\textendash 8}

Nine escalations are admitted, and not eleven. This build compares every
card with the first, where the naive verifier compared only the last, so it
catches two chains that the naive verifier missed. It is the repair a
careful reader might have tried, and it is still wrong.
Exercise~\ref{ex:anchor-build-two} finds the two chains. Four more builds
each remove one other guard: the signature check, the check that a card
names its parent, the comparison of expiry times, and the check of a card's
own expiry. For each, the kernel names a test that passes on the real core
and fails on that build.

The first two breaks and the fix are also ProVerif models, written for
this code (Listing~\ref{alg:anchor-build-models}).

\Needspace{16\baselineskip}
\pdrun[caption={ProVerif~2.05 on models of the card
check.},label={alg:anchor-build-models}]{models-cards}

A query of the form \texttt{not event(...)} that is reported \texttt{false}
means ProVerif found a run in which the event happens. So
\texttt{chain-root-only.pv} admits someone to \texttt{both}, which stands
for reading and writing, on a chain of three cards whose middle card grants
reading alone. \texttt{chain-per-hop.pv} admits no one that way. The third
query of each chain model is meant to be \texttt{false}: it shows that the
model can admit a chain at all, so the queries above it are not true only
because nothing can be admitted. In
\texttt{header-chooses-the-algorithm.pv} the verifier that reads the header
accepts a card that was never issued, and the verifier that does not read it
accepts none.

These models differ from v6 and v7 of \S\ref{sec:attn-soundness} in one way
that the chapter lists there as a gap. In \texttt{chain-per-hop.pv} the
verifier is given the harbor's key and no other. It reads each later key
from the parent card, as Listing~\ref{alg:anchor-build-admit} does, so these
models also cover which key a card must be signed with.
\pdscope{The count covers the 64 chains of three cards over two rights on
one path, with real signatures. The models cover the same sizes for any
number of sessions, with cryptography that cannot be broken. Longer chains
depend on the order being transitive (\S\ref{sec:delegation}). Tests check
that on a finite set of grants, and \S\ref{sec:implementation} on every
pattern of up to four bytes. No model here has two cards for one holder, so
the moved card is covered only by its run and its test.}

\begin{pdboundary}
\setlength{\parskip}{4pt plus 1pt}
A card in this build is still a bearer credential. Wren could not forge
one, and she did not need to: a card copied from Heron's files works for
whoever presents it until it expires, which in these runs is ten minutes
(\S\ref{sec:limitations}). \S\ref{sec:anchor-build-holder} builds the fix
for that. Nothing here revokes a card early
(\S\ref{sec:anchor-build-taken-back}). The verifier puts no limit on the length of a
chain. It trusts its signature library, \texttt{ed25519-dalek}, which no
check in this section examines. And \pdkernelname{} as a running service
does not yet require a card. The verifier is called by \texttt{harbor
scenario} and \texttt{harbor chains}, as above, and by the kernel's tests.
\end{pdboundary}
